\section*{Event Study}

% (h.1) Event date and trading windows
We now turn from annual return differences to the market reaction around a specific event. The event date is Friday, 24 June 2016, when the Brexit referendum result became known before trading. The $[-1,+1]$ window covers 23, 24 and 27 June, and $[-2,+2]$ covers 22, 23, 24, 27 and 28 June. Both exclude the intervening weekend.

% (h.2) Treatment definitions and sample composition
Using 2015 classifications, UK sales exposure means that at least 10\% of sales originate in the UK. The alternative treatment is a UK listing. The final panel contains 54 sales-exposed and 176 control firms, compared with 40 UK-listed and 190 control firms. Table~\ref{tab:exposure_groups} shows that the definitions disagree for 74 firms (32.174\%).

% Firm counts from code/04_task_h_event_study.R.
\begin{table}[H]
  \centering
  \singlespacing
  \caption{UK exposure classifications in the final sample}
  \label{tab:exposure_groups}
  \begin{tabular}{@{}lrrr@{}}
    \toprule
    UK sales share & Not UK-listed & UK-listed & Total \\
    \midrule
    Below 10\% & 146 & 30 & 176 \\
    At least 10\% & 44 & 10 & 54 \\
    Total & 190 & 40 & 230 \\
    \midrule
    \multicolumn{3}{@{}l}{Classified differently} & 74 \\
    \bottomrule
  \end{tabular}
\end{table}


\newpage
% (h.3) Tests based on the supplied group summaries
We use the supplied group-level average abnormal returns and standard errors. Cumulative abnormal returns sum daily abnormal returns. Their standard errors equal the daily standard error times $\sqrt{L}$ for $L$ trading days, assuming no correlation across days. For independent treated and control groups, $SE_{\mathrm{diff}}=\sqrt{SE_T^2+SE_C^2}$. Test statistics divide each return or difference by its standard error. Table~\ref{tab:event_returns} reports two-sided normal-approximation tests under both definitions.

% Maintained separately in LaTeX. Check values against the R analysis.
\begin{table}[H]
  \centering
  \singlespacing
  \caption{Abnormal returns around the Brexit referendum result}
  \label{tab:event_returns}
  \begin{tabular}{@{}llrrr@{}}
    \toprule
    Definition & Window & Treated & Control & Difference \\
    \midrule
    UK sales & Event day & $-4.678^{***}$ ($-28.682$) & $-0.032$ ($-0.371$) & $-4.646^{***}$ ($-25.108$) \\
     & $[-1,+1]$ & $-4.878^{***}$ ($-17.267$) & 0.097 (0.638) & $-4.975^{***}$ ($-15.522$) \\
     & $[-2,+2]$ & $-4.927^{***}$ ($-13.509$) & $-0.065$ ($-0.331$) & $-4.862^{***}$ ($-11.752$) \\
    \midrule
    UK listing & Event day & $-1.040^{***}$ ($-5.538$) & $-1.121^{***}$ ($-13.278$) & 0.081 (0.395) \\
     & $[-1,+1]$ & $-0.949^{***}$ ($-2.918$) & $-1.076^{***}$ ($-7.359$) & 0.127 (0.357) \\
     & $[-2,+2]$ & $-1.051^{**}$ ($-2.502$) & $-1.219^{***}$ ($-6.453$) & 0.168 (0.364) \\
    \bottomrule
  \end{tabular}

  \medskip
  \begin{minipage}{\textwidth}
    Notes. Group returns are in percent. Differences are treated minus control in percentage points. Test statistics are in parentheses. Event-day values are average abnormal returns. Longer windows report cumulative average abnormal returns. $^{*}p<0.10$, $^{**}p<0.05$, $^{***}p<0.01$.
  \end{minipage}
\end{table}


% (h.3--5) Results, robustness and interpretation
Sales-exposed firms lose 4.678\% on the event day, compared with 0.032\% for controls. The treated-minus-control difference is $-\,4.646$ percentage points and remains similar over three and five days ($-\,4.975$ and $-\,4.862$ points). All differences are significant at 1\%, whereas control-group returns are insignificant. Under the listing definition, both groups lose roughly 1\%, but none of their differences is significant.

UK sales capture exposure to UK demand and revenues, whereas a UK listing identifies a trading venue. The conclusion about differential losses therefore depends on the exposure measure. The supplied return file has no firm identifiers, so its aggregates cannot be recalculated for the filtered panel.
